% Public 3.1e interfaces; no new commands or options are used here.
\hitsetup{cauthor={测试作者},csubject={测试专业},csupervisor={测试导师},cdate={2026年10月}}
\cauthor{测试作者}
\ifdefined\HitBook
  \hitsetup{ctitle={旧接口测试},etitle={Legacy interface test},ckeywords={甲, 乙},ekeywords={alpha, beta}}
\else
  \ifdefined\HitArtplus
    \hitsetup{ctitleone={旧接口测试},ctitletwo={兼容验证}}
  \else
    \hitsetup{ctitlecover={旧接口测试}}
  \fi
\fi
\makeatletter
\ExplSyntaxOn
\cs_new_protected:Npn \HitProbeBool #1
  {\cs_if_exist:cT {ifhit@#1}
    {\use:c {ifhit@#1}
      \iow_term:n {HIT-PROBE:#1=true}
     \else
      \iow_term:n {HIT-PROBE:#1=false}
     \fi}}
\ExplSyntaxOff
\HitProbeBool{bachelor}\HitProbeBool{master}\HitProbeBool{doctor}\HitProbeBool{postdoc}
\HitProbeBool{harbin}\HitProbeBool{weihai}\HitProbeBool{shenzhen}
\HitProbeBool{english}\HitProbeBool{chinese}
\HitProbeBool{arialtoc}\HitProbeBool{arialtitle}\HitProbeBool{chapterbold}
\HitProbeBool{raggedbottom}\HitProbeBool{pifootnote}\HitProbeBool{glue}
\HitProbeBool{tocfour}\HitProbeBool{tocblank}\HitProbeBool{chapterhang}
\HitProbeBool{fulltime}\HitProbeBool{etitleauto}\HitProbeBool{subtitle}
\HitProbeBool{debug}\HitProbeBool{openright}\HitProbeBool{library}
\HitProbeBool{capcenterlast}\HitProbeBool{subcapcenterlast}\HitProbeBool{absupper}
\HitProbeBool{bsmainpagenumberline}\HitProbeBool{bsfrontpagenumberline}
\HitProbeBool{bsheadrule}\HitProbeBool{newtxmath}\HitProbeBool{splitbibitem}
\HitProbeBool{engtoc}\HitProbeBool{toc}\HitProbeBool{print}
\typeout{HIT-PROBE:cauthor=\hit@cauthor}
\typeout{HIT-PROBE:cdate=\hit@cdate}
\typeout{HIT-PROBE:textwidth=\the\textwidth}
\typeout{HIT-PROBE:textheight=\the\textheight}
\makeatother
\begin{document}
\ifdefined\HitCover
  \begin{cabstract}封面及选框回归。\end{cabstract}
  \begin{eabstract}Cover and checkbox regression.\end{eabstract}
  \makecover
\fi
\ifdefined\HitBook\mainmatter\chapter{兼容验证}[Compatibility check]\label{chapter}\fi
\section{正文}\label{section}
保留原版的命令、编号与排版。Legacy interfaces and layout remain available.
\ifdefined\HitBook
  \makeatletter 中文关键词：\hit@ckeywords。English keywords: \hit@ekeywords.\makeatother
\fi
\begin{figure}[htbp]\centering\rule{20mm}{4mm}\caption{图题}\label{figure}\end{figure}
\begin{figure}[htbp]\centering
  \subfigure[子图]{\rule{20mm}{3mm}\label{panel}}\caption{组合图}\label{panels}
\end{figure}
\begin{table}[htbp]\caption{表题}\label{table}\begin{tabular}{cc}甲&乙\\1&2\end{tabular}\end{table}
\begin{figure}[htbp]\centering
  \begin{minipage}{65mm}
    \rule{20mm}{3mm}
    \caption{多行标题用于验证最后一行居中算法及彩色文字的分组作用域，\textcolor{blue}{这里是蓝色文字}。}
    \label{long-figure}
  \end{minipage}
\end{figure}
\begin{longtable}{p{65mm}}
  \caption{长表格的多行标题用于验证共享的 caption 排版与编号接口，保持原版布局。}\label{long-table}\\
  表格正文。\\
\end{longtable}
\begin{equation}a=b\label{equation}\end{equation}
\begin{align}x&=y\label{align-a}\\u&=v\label{align-b}\end{align}
\ifdefined\HitTags
  \begin{equation}a=c\tag{S}\label{tagged}\end{equation}
  \begin{equation}b=d\tag*{T}\label{bare}\end{equation}
\fi
\begin{theorem}\label{theorem}定理正文。\end{theorem}
\begin{lemma}[测试]\label{lemma}引理正文。\end{lemma}
\begin{enumerate}\item 列表项。\label{item}\end{enumerate}
脚注\footnote{脚注正文。\label{note}}。
\ifdefined\HitStyle
  \begin{algorithm}[htbp]\caption{算法}\label{algorithm}$x\gets 1$\;\end{algorithm}
  \begin{lstlisting}[caption={Code},label={listing}]
int value = 1;
  \end{lstlisting}
\fi
\clearpage
图\ref{figure}，子图\ref{panel}，表\ref{table}，公式\eqref{equation}，
定理\ref{theorem}，引理\ref{lemma}，条目\ref{item}，脚注\ref{note}。
\ifdefined\HitTags\eqref{tagged}，\ref{bare}。\fi
\ifdefined\HitStyle\ref{algorithm}，\ref{listing}。\fi
\ifdefined\HitBook
  \appendix\chapter{附录}[Appendix]\label{appendix}
  \begin{equation}p=q\label{appendix-equation}\end{equation}
  \ref{appendix}，\eqref{appendix-equation}。
\fi
\end{document}
